\section{Statistical Considerations}

The statistical framework for this study is that of a confirmatory, randomized,
controlled superiority trial designed to test a comparative efficacy hypothesis
between the investigational regimen (Arm A) and institutional-standard care (Arm B).
Accordingly, the primary analysis is hypothesis-testing and is protected against type
I error inflation, in contrast to the predominantly descriptive framework of the
single-arm predicate \cite{kawchak2026phase1}. All analyses follow the principles of
the International Council for Harmonisation (ICH) E9 guideline on statistical
principles for clinical trials, including the estimand framework for defining the
target of estimation, are prespecified in this section and in the Statistical
Analysis Plan that this section summarizes, and are reported in accordance with the
CONSORT statement for parallel-group randomized trials \cite{Schulz2010consort} and
the TRIPOD+AI statement for the model-assisted advisory component
\cite{Collins2024tripodai}.

\subsection{Statistical Hypotheses}

The study has one primary efficacy hypothesis and a prespecified hierarchy of key
secondary hypotheses, each matched to a population-level summary and an analytic
method.

The primary efficacy hypothesis concerns progression-free survival (PFS). Let
$\theta$ denote the true hazard ratio (HR) for PFS comparing Arm A with Arm B. The
null and alternative hypotheses are $H_0\!: \theta \geq 1$ (no PFS benefit of Arm A
relative to Arm B) versus $H_1\!: \theta < 1$. The hypothesis is tested by a
stratified log-rank test using the randomization strata (resectability, KRAS allele,
neoadjuvant therapy, and site), and the treatment effect is estimated by a stratified
Cox proportional-hazards model that yields the HR with a two-sided $95\%$ confidence
interval. The test is two-sided at the $\alpha = 0.05$ level, with the type I error
allocated across one interim and the final analysis by a group-sequential alpha
spending function ($\S$9.2, $\S$9.4.6).

The key secondary hypotheses are tested only if the primary PFS hypothesis is
rejected, in a fixed-sequence (hierarchical) order that controls the family-wise type
I error at two-sided $0.05$ without further multiplicity adjustment: overall survival
(OS); the margin-negative (R0) resection rate; the clinically relevant (ISGPS Grade B
or C) postoperative pancreatic fistula rate \cite{Bassi2017}; the major pathologic
response (MPR) rate; and the week-12 KRAS ctDNA clearance rate \cite{Tie2019ctdna}.
Each is tested at two-sided $0.05$, and testing proceeds down the sequence only while
each successive null hypothesis is rejected; the first non-rejection stops the
confirmatory sequence, and endpoints below that point are reported descriptively. The
remaining secondary endpoints are estimation targets reported with two-sided $95\%$
confidence intervals and are not part of the confirmatory hierarchy. ICH E9 estimand
attributes (population, endpoint, handling of intercurrent events such as open
conversion, withdrawal, new anticancer therapy, or death, and the population-level
summary) are specified for every endpoint in the Statistical Analysis Plan.

\subsection{Sample Size Determination}

The planned sample size is $n=220$ participants randomized 1:1 (110 per arm). The
sample size is event-driven: under proportional hazards, approximately $140$ PFS
events provide $85\%$ power at a two-sided $\alpha = 0.05$ to reject the primary null
hypothesis when the true PFS hazard ratio is $0.60$, corresponding to a median PFS of
$14.0$ months in Arm B and $23.3$ months in Arm A. The randomization of 220
participants is projected to accrue the required $140$ events over approximately a
30-month accrual period across the eight sites followed by a minimum of 18 months of
additional follow-up. To preserve the evaluable sample against non-evaluability
(major eligibility violations, withdrawal before any post-baseline assessment, or
missing the primary assessment), up to approximately $245$ participants may be
enrolled, allowing for up to $10\%$ non-evaluability while retaining $220$ evaluable
randomized participants.

A single interim analysis for efficacy and futility is performed when approximately
$60\%$ of the targeted PFS events (approximately $84$ events) have accrued. The type
I error is controlled across the interim and final analyses by the Lan-DeMets alpha
spending function approximating the O'Brien-Fleming boundary, which spends little
alpha at the interim and preserves close to the full two-sided $0.05$ for the final
analysis \cite{DeMets1994lan,Jennison2000group}; a non-binding futility boundary
based on the corresponding beta spending is evaluated at the same interim. The
operating characteristics of the design are summarized in Table~\ref{tab:power}.

\begin{table}[htbp]
\centering
\footnotesize
\caption{Design and operating characteristics for the confirmatory
progression-free-survival comparison. The trial is an event-driven, randomized,
two-arm superiority design with one group-sequential interim analysis governed by a
Lan-DeMets O'Brien-Fleming alpha spending function.}
\label{tab:power}
\begin{tabularx}{\textwidth}{L{3cm} L{3.2cm} Y}
\toprule
\textbf{Parameter} & \textbf{Planned value} & \textbf{Basis and operating characteristic} \\
\midrule
Design & Randomized 1:1, two-arm, superiority & Multicenter (8 sites), controlled, open-label with blinded independent central review. \\
Arms & Arm A vs Arm B & Arm A: daraxonrasib RP2D plus LLM-directed robotic Whipple. Arm B: modified FOLFIRINOX plus standard pancreaticoduodenectomy \cite{Conroy2018}. \\
Allocation & 1:1 (110 per arm) & Central permuted-block randomization, stratified by resectability, KRAS allele, neoadjuvant therapy, and site. \\
Primary endpoint & Progression-free survival & Time from randomization to RECIST 1.1 progression or death; stratified log-rank with stratified Cox HR. \\
Target hazard ratio & HR $= 0.60$ & Alternative hypothesis ($H_1\!: \theta < 1$); corresponds to the assumed treatment benefit. \\
Median PFS & 14.0 vs 23.3 months & Arm B (control) versus Arm A (experimental) under proportional hazards. \\
Power & $85\%$ & At HR $= 0.60$ and two-sided $\alpha = 0.05$. \\
Significance level & $\alpha = 0.05$ (two-sided) & Allocated across one interim and the final analysis by alpha spending. \\
Target events & approximately 140 PFS events & Event-driven primary analysis; database lock at the target event count. \\
Interim analysis & one, at approximately $60\%$ of events ($\sim 84$) & Efficacy and futility; Lan-DeMets O'Brien-Fleming spending; DSMB review \cite{DeMets1994lan,Jennison2000group}. \\
Sample size & $n = 220$ (up to $\sim 245$ enrolled) & Up to $\sim 245$ enrolled to allow $\leq 10\%$ non-evaluability. \\
Accrual and follow-up & $\sim 30$-month accrual, $\geq 18$-month follow-up & Across 8 sites; follow-up to the 24-month overall-survival endpoint. \\
\bottomrule
\end{tabularx}
\end{table}

\subsection{Populations for Analyses}

Four analysis populations are defined and applied consistently across all analyses.
The intention-to-treat (ITT) population comprises all randomized participants
analyzed according to the arm to which they were randomized, and is the primary
population for the primary and key secondary efficacy analyses. The modified
intention-to-treat (mITT) population comprises all randomized participants who
received any study intervention and contributed at least one post-baseline
assessment, and supports sensitivity analyses. The per-protocol population comprises
participants who received the assigned intervention without a major protocol deviation
that could affect the efficacy assessment, and supports supportive analyses of the
primary and surgical-quality endpoints. The safety population comprises all
participants who received any study intervention, analyzed according to the
intervention actually received, and is the population for all safety analyses. The
derivation of the populations and the group-sequential interim with the DSMB efficacy
and futility boundaries are shown in Figure 18.

\begin{mermaidfig}
\node[mmin]  (rand) at (0.3,-1.15)  {Randomized 1:1\\n = 220\\(110 per arm)};
\node[mmstep](itt)  at (4.4,2.3)    {ITT population\\all randomized\\(primary)};
\node[mmstep](mitt) at (4.4,0)      {modified ITT\\dosed with $\geq 1$\\post-baseline assessment};
\node[mmstep](pp)   at (4.4,-2.3)   {Per-protocol\\assigned intervention,\\no major deviation};
\node[mmstep](saf)  at (4.4,-4.6)   {Safety population\\intervention\\as received};
\node[mmdec] (intm) at (10.3,2.3)   {Interim at $\sim 60\%$\\events ($\sim 84$):\\DSMB review};
\node[mmgoal](bound)at (8,-0.6)     {O'Brien-Fleming\\efficacy / futility\\boundary crossed};
\node[mmgoal](cont) at (11.7,-0.215){Continue to\\final analysis};
\draw[mmedge] (rand.north east) -- (itt.south west);
\draw[mmedge] (rand.north east) -- (mitt.west);
\draw[mmedge] (rand.south east) -- (pp.west);
\draw[mmedge] (rand.south east) -- (saf.north west);
\draw[mmedge] (itt) -- (intm);
\draw[mmedge] (intm.south west) -- node[left,font=\scriptsize, yshift=0.2cm, xshift=0.1cm]{crossed} (bound.north);
\draw[mmedge] (intm.south east) -- node[left,font=\scriptsize,pos=0.6, yshift=0.1cm, xshift=-0.1cm]{within} (cont.north);
\end{mermaidfig}
\vspace{-0.6cm}
\figcaption{Figure 18. The four analysis populations (intention-to-treat,
modified intention-to-treat, per-protocol, safety) and the single group-sequential
interim analysis at approximately $60\%$ of the targeted progression-free-survival
events, at which the DSMB evaluates the Lan-DeMets O'Brien-Fleming efficacy and
futility boundaries. Binding device-safety halt rules operate in parallel: a
device-related serious-adverse-event rate breach or a hard-limit breach triggers a
mandatory DSMB review.}

\subsection{Statistical Analyses}

\subsubsection{General Approach}

The primary and key secondary efficacy analyses are confirmatory and are protected
against type I error inflation; the remaining analyses are estimation-based and
descriptive. Time-to-event endpoints are analyzed by the stratified log-rank test and
summarized by the Kaplan-Meier method with the number at risk, and the treatment
effect is estimated by a stratified Cox proportional-hazards model reporting the
hazard ratio with a two-sided $95\%$ confidence interval; the proportional-hazards
assumption is assessed and, if it is not supported, a prespecified alternative
(restricted mean survival time) is reported. Binary endpoints are summarized as
proportions with two-sided $95\%$ confidence intervals and compared by a
stratified test (Cochran-Mantel-Haenszel) where a confirmatory comparison is
specified. Continuous variables are summarized by the number of non-missing
observations, mean, standard deviation, median, minimum, and maximum; categorical
variables are summarized by counts and percentages. ICH E9 estimands govern the
handling of intercurrent events. Missing data are handled by the prespecified primary
estimand strategy, with sensitivity analyses to assess robustness. All analyses are
performed in a validated, version-controlled statistical environment under a
deterministic random seed, and the analysis code and derived data sets are deposited
so that every reported result is reproducible.

\subsubsection{Analysis of the Primary Endpoint}

The primary endpoint, PFS, is analyzed on the ITT population. The primary analysis is
the stratified log-rank test of $H_0\!: \theta \geq 1$ against $H_1\!: \theta < 1$ at
the alpha allocated to the final analysis by the group-sequential spending function,
with the hazard ratio and its two-sided $95\%$ confidence interval estimated from the
stratified Cox model. The primary read is the investigator-assessed PFS; a
prespecified sensitivity analysis repeats the test and estimation on the
BICR-assessed PFS to confirm that the open-label delivery does not propagate into
endpoint ascertainment ($\S$8.1). Additional sensitivity analyses repeat the primary
analysis on the per-protocol and mITT populations and under alternative
censoring rules for new anticancer therapy.

\subsubsection{Analysis of the Key Secondary Endpoints}

The key secondary endpoints are analyzed only if the primary null hypothesis is
rejected, in the fixed-sequence hierarchy of $\S$9.1: OS, then the R0 resection rate,
then the ISGPS Grade B or C fistula rate, then the MPR rate, then the week-12 KRAS
ctDNA clearance rate. OS is analyzed by the stratified log-rank test and the
stratified Cox model on the ITT population; the binary key secondary endpoints are
analyzed by the stratified Cochran-Mantel-Haenszel test with the risk difference and
its $95\%$ confidence interval. Each test is conducted at two-sided $0.05$, and the
sequence stops at the first non-rejection. The key secondary and the remaining
secondary endpoints, their analysis populations and methods, and the external Dutch
2025 nationwide robotic-pancreatoduodenectomy descriptive context are summarized in
Table~\ref{tab:secendpts} \cite{DutchCohort2025}.

\begin{table}[htbp]
\centering
\footnotesize
\caption{Key secondary (confirmatory, fixed-sequence) and other secondary
(estimation) endpoints, the analysis population and method, and the reporting and
external comparator. The Dutch 2025 nationwide robotic-pancreatoduodenectomy cohort
is an external descriptive reference, not an inferential comparator.}
\label{tab:secendpts}
\begin{tabularx}{\textwidth}{L{4.5cm} L{3.4cm} Y}
\toprule
\textbf{Endpoint} & \textbf{Population / method} & \textbf{Reporting and comparator} \\
\midrule
Overall survival (key) & ITT; stratified log-rank, Cox HR & First in the hierarchy after PFS; KM curves, HR with $95\%$ CI; 24-month rate. \\
R0 resection rate (key) & ITT; stratified CMH & Risk difference with $95\%$ CI; central masked synoptic pathology. \\
ISGPS Grade B/C fistula (key) & ITT; stratified CMH & Risk difference with $95\%$ CI; graded per the 2016 ISGPS update \cite{Bassi2017}. \\
Major pathologic response (key) & ITT; stratified CMH & Risk difference with $95\%$ CI; central pathology grading. \\
Week-12 KRAS ctDNA clearance (key) & ITT; stratified CMH & Risk difference with $95\%$ CI; central laboratory assay \cite{Tie2019ctdna}. \\
Clavien-Dindo grade III+ & Safety; proportion, $95\%$ CI & Estimation only; Dutch 2025 reference 41.3\% \cite{DutchCohort2025}. \\
90-day mortality & Safety; proportion, $95\%$ CI & Estimation only; Dutch 2025 reference 3.9\% \cite{DutchCohort2025}. \\
Conversion to open & Safety; proportion, $95\%$ CI & Estimation only; Dutch 2025 reference 10.1\% \cite{DutchCohort2025}. \\
Operative time, EBL, length of stay & Per-protocol; descriptive & Mean / median with range; Dutch 2025 descriptive context \cite{DutchCohort2025}. \\
Quality of life & mITT; descriptive / mixed model & EORTC QLQ-C30 and PAN26 over time \cite{EORTCqlqc30}. \\
12- and 18-month PFS rate; 24-month OS rate & ITT; Kaplan-Meier & Landmark rates with $95\%$ CI. \\
\bottomrule
\end{tabularx}
\end{table}

\subsubsection{Safety Analyses}

All safety analyses are performed on the safety population, with events analyzed
according to the intervention actually received. Adverse events are coded to the
Medical Dictionary for Regulatory Activities (MedDRA) and tabulated by system organ
class and preferred term, by maximum severity using the Common Terminology Criteria
for Adverse Events (CTCAE), and by relationship to the drug, to the device (Arm A),
and to the procedure, with proportions accompanied by $95\%$ confidence intervals and
reported by arm. Serious adverse events and deaths are listed individually. Device
telemetry constitutes a second, parallel safety stream unique to Arm A. The
cyber-physical record is summarized per procedure and in aggregate: emergency-stop
activations and the measured stop latency against the $\leq 500$ ms requirement;
positional and force excursions against the per-arm $\leq 3$ N and cumulative
$\leq 18$ N caps; vascular safety-zone gating events; and the counts of LLM
advisories that were accepted, blocked, or escalated to the human operator. A
dedicated Physical AI adverse-event tabulation records, for every event, whether an
LLM advisory, a robot command, a safety-limit trip, or a human override was
implicated, so that events arising from the autonomy-bearing layer are visible and
traceable to a deterministic seed and a repository commit through the federated
hash-chained audit trail ($\S$8.3.6). Laboratory data, vital signs, and
physical-examination findings are summarized descriptively with shift tables relative
to baseline.

\subsubsection{Baseline Descriptive Statistics}

Baseline characteristics are summarized descriptively by arm and overall. Reported
variables include demographics (age, sex, race, ethnicity); disease staging and
resectability category (resectable versus borderline-resectable); the specific KRAS
G12 mutation; baseline CA 19-9; Eastern Cooperative Oncology Group (ECOG) performance
status; prior neoadjuvant therapy and number of cycles; and relevant comorbidity and
organ-function indices. Consistent with CONSORT and the randomized design, baseline
characteristics are presented by arm without a formal significance test of baseline
balance \cite{Schulz2010consort}; any notable imbalance in a strong prognostic factor
is addressed by a prespecified covariate-adjusted sensitivity analysis of the primary
endpoint.

\subsubsection{Planned Interim Analyses}

A single formal interim analysis of the primary endpoint is conducted when
approximately $60\%$ of the targeted PFS events (approximately $84$ events) have
accrued, and is reviewed by the independent DSMB. Efficacy is assessed against the
Lan-DeMets O'Brien-Fleming efficacy boundary and futility against the corresponding
non-binding futility boundary, so that the cumulative two-sided type I error across
the interim and final analyses is controlled at $0.05$
\cite{DeMets1994lan,Jennison2000group}. In parallel with the efficacy interim, binding
device-safety halt rules operate throughout the trial: a device-related serious
adverse-event rate that reaches the prespecified halt threshold, or a breach of a hard
cyber-physical safety limit (emergency-stop latency, or the per-arm or cumulative
force caps), triggers a mandatory DSMB review and possible halt, and a halt for an
immediate hazard is implemented at once under $\S$312.30(h). No efficacy stopping
decision is taken outside the prespecified group-sequential boundary.

\subsubsection{Sub-Group Analyses}

Prespecified subgroup analyses of the primary endpoint are conducted by resectability
category, KRAS allele (G12D versus other G12), neoadjuvant therapy (yes versus no),
and site, with the hazard ratio and $95\%$ confidence interval displayed in a forest
plot and a test of treatment-by-subgroup interaction reported. These analyses are
prespecified but are interpreted with caution as supportive and hypothesis-refining;
they are not adjusted for multiplicity, and no confirmatory conclusion is drawn from
any single subgroup.

\subsubsection{Tabulation of Individual Participant Data}

Individual participant data are tabulated to support the aggregate analyses.
Participant-level listings are provided for demographics and baseline characteristics;
randomized arm, stratification factors, and actual intervention received; all adverse
events with onset, severity, seriousness, relationship, and outcome; serious adverse
events and deaths; the efficacy endpoints and their event or assessment dates; and,
for Arm A, the per-procedure device-telemetry and Physical AI event summaries, each
keyed to participant and timepoint. Listings are presented so that every aggregate
statistic can be traced to its contributing records.

\subsubsection{Exploratory Analyses}

Exploratory analyses are non-confirmatory and are clearly labeled as such. They
include the concordance between LLM advisories and the verdicts of the deterministic
safety gates; the sim-to-real gap between the Phase 0 simulation predictions and the
observed intraoperative telemetry; the performance of the federated-learning model
across sites over the course of the trial; the longitudinal dynamics of KRAS ctDNA
beyond the week-12 clearance endpoint \cite{Tie2019ctdna}; and a health-economic and
cost-effectiveness analysis linked to the co-investment facility. These exploratory
results generate hypotheses for later development and are not used to draw confirmatory
efficacy conclusions in this protocol.
